\section{Goals must refer to external state}
\label{sec:C10:goals}
A plan is a policy for changing or learning about a task state, not merely an ordered explanation. Let $G(E)$ be a goal predicate on environment state, $I(S)$ an invariant on complete system state, and $c(a,S)$ an action cost. A useful planning problem specifies which terminal states satisfy the user objective, which intermediate states are prohibited, and what resources are available. An agent that produces a convincing account of $G$ without establishing it in the environment has not completed that planning problem. The distinction is especially important when a language model can simulate the expected wording of a successful tool response.

Goal predicates can be partial or ambiguous. Passing tests may be an operational acceptance condition for a coding task without fully specifying intended behavior. A reservation can exist in a system while violating a requested time or budget. The planner should retain unresolved requirements and seek clarification or evidence when those distinctions affect action. Converting an ambiguous instruction into one confident internal goal silently changes the task. A practical goal representation therefore includes source instructions, interpreted constraints, acceptance checks, and unresolved assumptions, rather than only a short objective summary.

Multiobjective planning requires explicit priorities. Success, latency, monetary cost, privacy, and permission compliance cannot always be compressed into one unconstrained scalar reward. Some are hard feasibility conditions; others are tradeoffs within the feasible set. A planner should not spend an unauthorized action to gain a slightly higher task score. Let $Y_I(\tau)\in\{0,1\}$ indicate whether trajectory $\tau$ contains a state violating the declared invariant $I$. A constrained objective can require $P_\pi(Y_I=1)\leq\delta$, in addition to statewise authorization. These formulations are meaningful only when the probabilities and predicates have an evidential basis. A model-generated confidence statement does not establish the risk bound.

The plan's initial state is also a claim requiring evidence. A resumed session can contain stale memory, incomplete operations, or a tool schema different from the previous run. Before continuing, reconcile the durable event history with the current environment and permissions. The policy may need an information-gathering action rather than an execution action. This makes planning inseparable from context and memory: the planner's apparent reasoning competence can be undermined by an incorrect state estimate before the first proposed step.

